% Part: lambda-calculus
% Chapter: introduction
% Section: unique-readability

\documentclass[../../../include/open-logic-section]{subfiles}

\begin{document}

\olfileid{lam}{syn}{unq}
\olsection{একক পাঠযোগ্যতা}

প্রতিটি পদ গঠনের উপায় একক কি না, সেই প্রশ্ন উঠতে পারে; উত্তরটি হ্যাঁ।
প্রতিটি ল্যাম্বডা পদ নির্মাণ ও ব্যাখ্যা করার উপায় একটিই। নিচের আলোচনায়
\emph{গঠনপ্রণালী} বলতে \olref[trm]{defn:term}-এর গঠন-নিয়ম একবার বা
একাধিকবার ব্যবহার করে কোনো পদ নির্মাণের প্রক্রিয়াকে বোঝানো হবে।

\begin{lem}\ollabel{lem:term-start}
কোনো পদ হয় একটি চলরাশি, নয়তো একটি বন্ধনী দিয়ে শুরু হয়।
\end{lem}

\begin{proof}
কোনো কিছু কেবল \olref[trm]{defn:term} অনুসারে নির্মিত হলেই পদ হিসেবে
গণ্য হয়। সেটি \olref[trm]{defn:term-var}-এর ফল হলে অবশ্যই একটি চলরাশি।
আর \olref[trm]{defn:term-abs} বা \olref[trm]{defn:term-app}-এর ফল হলে
সেটি একটি বন্ধনী দিয়ে শুরু হয়।
\end{proof}

\begin{lem}\ollabel{lem:app-start}
কোনো প্রয়োগের ফল হয় পরপর দুটি বন্ধনী, নয়তো একটি বন্ধনী ও একটি চলরাশি
দিয়ে শুরু হয়।
\end{lem}

\begin{proof}
$M$ কোনো প্রয়োগের ফল হলে সেটি $(PQ)$ রূপের, তাই একটি বন্ধনী দিয়ে শুরু
হয়। $P$ একটি পদ বলে \olref{lem:term-start} অনুসারে সেটি হয় একটি বন্ধনী,
নয়তো একটি চলরাশি দিয়ে শুরু হয়।
\end{proof}

\begin{lem}\ollabel{lem:initial}
কোনো পদের প্রকৃত পূর্বাংশ নিজে একটি পদ নয়।
\end{lem}

\begin{prob}
পদের দৈর্ঘ্যের উপর আরোহ প্রয়োগ করে \olref[lam][syn][unq]{lem:initial}
প্রমাণ করো।
\end{prob}

\begin{prop}[একক পাঠযোগ্যতা] \ollabel{prop:unq}
প্রতিটি পদের গঠনপ্রণালী একক। অন্যভাবে বললে, কোনো পদ~$M$ একটি
গঠনপ্রণালীতে গঠিত হলে ওই পদটি গঠন করতে পারে এমন গঠনপ্রণালী সেটিই একমাত্র।
\end{prop}

\begin{proof}
  পদের গঠনের উপর আরোহ করে এটি প্রমাণ করি।
  \begin{enumerate}
    \item $M$ হলো কোনো চলরাশি~$x$-এর রূপের, যেখানে $x$ একটি চলরাশি। বিমূর্তন ও প্রয়োগের ফল
      সবসময় বন্ধনী দিয়ে শুরু হয় বলে সেগুলি দিয়ে~$M$ নির্মিত হতে পারে
      না। অতএব $M$-এর গঠনপ্রণালী অবশ্যই
      \olref[trm]{defn:term}\olref[trm]{defn:term-var}-এর একটিমাত্র ধাপ।
    \item $M$ হলো $(\lambd[x][N])$ রূপের, যেখানে $x$ কোনো চলরাশি এবং
      $N$ একটি পদ। এটি একটিমাত্র চলরাশি নয় বলে
      \olref[trm]{defn:term}\olref[trm]{defn:term-var} অনুসারে নির্মিত হতে
      পারে না। \olref{lem:app-start} অনুসারে এটি কোনো প্রয়োগের ফলও নয়।
      তাই $M$ কেবল~$N$-এর উপর বিমূর্তনের ফল হতে পারে। আরোহানুমান থেকে
      জানি, $N$-এর গঠনপ্রণালী নিজেই একক।
    \item $M$ হলো $(PQ)$ রূপের, যেখানে $P$ ও $Q$ পদ। এটি বন্ধনী দিয়ে
      শুরু বলে \olref[trm]{defn:term}\olref[trm]{defn:term-var} অনুসারে
      নির্মিত হতে পারে না। \olref{lem:term-start} অনুসারে $P$ কখনো
      $\lambd$ দিয়ে শুরু হতে পারে না, তাই $(PQ)$ কোনো বিমূর্তনের ফলও হতে
      পারে না। এবার ধরি, প্রয়োগের মাধ্যমে $M$ নির্মাণের আরেকটি উপায় আছে;
      যেমন, সেটি $(P'Q')$ রূপেরও। তাহলে $P$, $P'$-এর প্রকৃত পূর্বাংশ
      (অথবা উল্টোটি); কিন্তু \olref{lem:initial} অনুসারে তা অসম্ভব। অতএব
      $P$ ও~$Q$ এককভাবে নির্ধারিত, এবং আরোহানুমান থেকে $P$ ও~$Q$-এর
      গঠনপ্রণালীও একক।
  \end{enumerate}
\end{proof}

উপরের প্রস্তাবটির আরও পাঠযোগ্য ভাষান্তর নিচে দেওয়া হলো:
\begin{prop}
  কোনো পদ~$M$ কেবল নিচের রূপগুলির একটির হতে পারে:
  \begin{enumerate}
    \item $x$, যেখানে $x$ হলো~$M$ দ্বারা এককভাবে নির্ধারিত একটি চলরাশি।
    \item $(\lambd[x][N])$, যেখানে $x$ একটি চলরাশি এবং $N$ আরেকটি পদ;
      দুটিই~$M$ দ্বারা এককভাবে নির্ধারিত।
    \item $(PQ)$, যেখানে $P$ ও~$Q$ হলো~$M$ দ্বারা এককভাবে নির্ধারিত দুটি
      পদ।
  \end{enumerate}
\end{prop}

\end{document}
